\chapter{Problemi numerici e stabilità}


\section{Problemi numerici}

Un problema numerico è un problema matematico in cui i dati sono espressi mediante numeri 
(o vettori di numeri) e la soluzione cercata è anch'essa un numero (o un insieme finito di numeri)
che soddisfa una certa relazione o condizione specificata, come per esempio un'equazione,
una condizione di equilibrio, un sistema di equazioni, una condizione di massimo o minimo, oppure una proprietà di congruenza o proporzionalità.

\begin{defi}[Problema Diretto]
Il problema diretto consiste nel calcolare la \textcolor{blue}{soluzione} (\textcolor{blue}{$x$}) di una funzione $F$ dato un \textcolor{blue}{set di dati} specifico (\textcolor{blue}{$d \epsilon D$}), in simboli:
\[
\textcolor{blue}{trova \, {x} \mid F(x, d) = 0 \quad x,d \, \epsilon \, \mathbb{R}^n}
\]

\end{defi}

\section{Problemi stabili e malcondizionati}

Esistono problemi tali che, qualunque algoritmo venga utilizzato per risolverli,
l'errore generato nel risultato risulta elevato e tale da rendere privo di significato il risultato stesso. \\
Per questi problemi, detti $malcondizionati$, piccole variazioni nei dati inducono grosse variazioni nei risultati.

\begin{defi}[Stabilità]
Si consideri un generico problema diretto. Sia \textcolor{blue}{$\delta_d$} una \textcolor{blue}{perturbazione nei dati} e \textcolor{blue}{$\delta_x$}
la corrispondente variazione nella soluzione, tale che:
\[
\textcolor{blue}{F(x+\delta_x, d+\delta_d)=0}
\] 
allora il problema è detto \textcolor{blue}{stabile} se:
\[
\textcolor{blue}{\exists k_0(d) \in \mathbb{R}^n : \|\delta_x\| \leq k_0(d) \|\delta_d\|}
\]
dove \textcolor{blue}{$k_0(d)$} è una costante che non dipende da \textcolor{blue}{$\delta_d$}
o \textcolor{blue}{$\delta_x$}.

Si parla di \textcolor{blue}{stabilità relativa} se:
\[
\textcolor{blue}{\exists k_0(d) \in \mathbb{R}^n : \frac{\|\delta_x\|}{\|x\|} \leq k_0(d) \frac{\|\delta_d\|}{\|d\|}}
\]
\end{defi}


\subsection{Numero di Condizionamento}

Definiamo ora il concetto di \textcolor{blue}{numero di condizionamento}.

Questo, altro non è che una misura della sensibilità del problema alle perturbazioni nei dati.

\begin{defi}[Numero di Condizionamento]
    Si consideri un generico problema diretto con \textcolor{blue}{$x\neq0$} 
    e \textcolor{blue}{$d\neq0$}. Definiamo:
    \begin{itemize}
        \item \textcolor{blue}{$K_{abs}(d)$}: numero di condizionamento assoluto del problema.\[\textcolor{blue}{K_{abs}(d) = \sup\left\{ \frac{\|\delta_x\|}{\|\delta_d\|}; \delta_d \neq 0\right\}} \]
        \item \textcolor{blue}{$K(d)$}: numero di condizionamento relativo del problema.\[\textcolor{blue}{K(d) = \sup\left\{\frac{\|\delta_x\|}{\|x\|}\frac{\|d\|}{\|\delta_d\|}; \delta_d \neq 0\right\}}\]
    \end{itemize}
\end{defi}

\begin{exercise}
Esempio di calcolo del numero di condizionamento per un problema specifico:
Supponiamo di avere un problema diretto definito da una funzione $F(x, d) = 0$ dove:
\[
 F(x, d) = x - (d_1+d_2) = 0
\]
Adottando sul nostro spazio la norma $\|\cdot\|_{l_1}$ dimostra che il numero di condizionamento relativo del problema è $K(d) = 1$.
\end{exercise}

\begin{obs}
    La norma $\|\cdot\|_{l_1}$ nell'esempio agisce sia su $x \epsilon \mathbb{R}$ che su $d \epsilon \mathbb{R}^2$.
\end{obs}

Se la soluzione $(x)$ di un generico problema diretto esiste ed è unica, allora: 
\[
\exists ! G(d) = x  \qquad \text{tale che} \,\, F(G(d), d) = 0
\]
Dove $G$ è un'applicazione detta \textcolor{blue}{risolvente} del problema diretto.

In questo caso esiste una relazione diretta tra il numero di condizionamento del problema e la derivata della risolvente.

\begin{theorem}
    Dato un generico problema diretto con soluzione unica e risolvente $G$ supposta \textcolor{blue}{continua}
    e \textcolor{blue}{derivabile}, allora:
    \[\textcolor{blue}{K(d) \leq \|G'(d)\| \frac{\|d\|}{\|G(d)\|}}\]
    Il numero di condizionamento assoluto del problema è invece dato da:
    \[\textcolor{blue}{K_{abs}(d) \approx \|G'(d)\|}\]
\end{theorem}

\begin{obs}
Il teorema precedente fornisce un criterio pratico per stimare il numero di condizionamento di un problema diretto a partire dalla derivata della sua risolvente, che nella maggior parte dei casi risulta più semplice da calcolare o approssimare rispetto al numero di condizionamento stesso.

Questa relazione è anche coerente con l'intuizione: la derivata misura proprio la sensibilità della soluzione rispetto a piccole variazioni dei dati, ed è quindi naturale che sia strettamente legata al numero di condizionamento del problema.
\end{obs}

\begin{proof}
    Supponendo che $G$ sia continua e derivabile si ha:
    \[
    x+\delta_x = G(d+\delta_d) = G(d) + G'(d)\delta_d + o(\|\delta_d\|); \quad \delta_d \to 0
    \]
    da cui:
    \[\delta_x = G(d+\delta_d) - x = G(d+\delta_d) - G(d) \approx G'(d)\delta_d + o(\|\delta_d\|)\]
    e quindi:
    \[
    k(d) = \frac{\|\delta_x\|}{\|x\|} \frac{\|d\|}{\|\delta_d\|} \approx \frac{\|G'(d)\delta_d\|}{\|x\|}\frac{\|d\|}{\|\delta_d\|} \leq \|G'(d)\| \frac{\|d\|}{\|x\|}
    \]
    Dove nell'ultimo passaggio abbiamo utilizzato la disuguaglianza di Schwartz della norma.

    L'approssimazione del numero di condizionamento assoluto segue da un ragionamento analogo.  
\end{proof}

\section{Metodi numerici}
Come abbiamo già osservato, non è sempre possibile risolvere analiticamente un problema matematico.

Per questo motivo, introduciamo il concetto di \textcolor{blue}{metodo numerico}.

\begin{defi}[Metodo numerico]
    Un \textcolor{blue}{metodo numerico} è una procedura di calcolo che serve a trovare una soluzione approssimata di un problema matematico,
    soprattutto quando la soluzione esatta è difficile o impossibile da ottenere in forma analitica.

    Un metodo numerico consiste in una \textcolor{blue}{successione} di \textcolor{blue}{problemi approssimati}:
    \[
        \textcolor{blue}{(1)_n \quad F_n(x_n, d_n) = 0}
    \]
    con \textcolor{blue}{$d_n \to d$} e \textcolor{blue}{$F_n$} \textit{approssima} \textcolor{blue}{$F$} quando $n \to \infty$.
\end{defi}

La nostra speranza è che, risolvendo i problemi approssimati, si riesca ad ottenere una successione di soluzioni approssimate \textcolor{blue}{$x_n$} che converga alla soluzione esatta \textcolor{blue}{$x$} del problema diretto originale.

\begin{defi}[Convergenza di un metodo numerico]
    Un metodo numerico è detto \textcolor{blue}{convergente} se l'errore \textcolor{blue}{$e_n$} associato al metodo
    tende a zero quando $n$ tende all'infinito, ovvero se:
    \[
        \textcolor{blue}{\lim_{n \to \infty} e_n = \lim_{n \to \infty} \|x_n - x\| = 0}
    \]
\end{defi}

\noindent\textcolor{red}{\faWarning\ \textbf{Attenzione.}} Da qui in poi supporremo di avere sempre a che fare con problemi diretti stabili e metodi numerici convergenti, altrimenti i risultati che otterremo potrebbero non essere validi.

\begin{defi}[Consistenza di un metodo numerico]
    Un metodo numerico è detto \textcolor{blue}{consistente} se:
    \[
        \textcolor{blue}{\lim_{n \to \infty} F_n(x, d) = F(x, d)}
    \]
    per ogni \textcolor{blue}{$x$} e \textcolor{blue}{$d \in \mathbb{D}_n$}.
\end{defi}

\begin{obs}
    La consistenza di un metodo numerico ci assicura che l'errore di approssimazione tra l'operatore continuo originale $F$ e l'operatore discreto e approssimato $F_n$ tende a zero quando $n$ tende all'infinito:
    \[
        \textcolor{blue}{\lim_{n \to \infty} \|F_n(x, d) - F(x, d)\| = \|F_n(x, d)\| = 0}
    \]
\end{obs}

In altre parole, la consistenza ci garantisce che il metodo numerico stia effettivamente approssimando il problema originale in modo sempre più accurato.

\begin{defi}[Metodo fortemente consistente]
    Un metodo numerico è detto \textcolor{blue}{fortemente consistente} se:
    \[
        \textcolor{blue}{\|F_n(x, d) - F(x, d)\| = O(\|x_n - x\|)}
    \]
\end{defi}

\begin{obs}[Idea intuitiva]
    La quantità
    \[
        r_n(x,d) := F_n(x,d)-F(x,d)
    \]
    misura il \textcolor{blue}{difetto di modello} introdotto dalla discretizzazione.

    Dire che il metodo è fortemente consistente significa che questo difetto non è ``più grande'' dell'errore sulla soluzione:
    \[
        \|r_n(x,d)\| \le C\,\|x_n-x\| \qquad (n \text{ grande})
    \]
    per una costante $C>0$ indipendente da $n$.

    In parole semplici: quando $x_n$ si avvicina a $x$, anche l'errore di consistenza si riduce almeno con lo stesso ordine. Quindi la discretizzazione è sufficientemente fedele al problema continuo e non ``rovina'' la qualità dell'approssimazione numerica.
\end{obs}

Introduciamo ora un teorema il cui contenuto è intuitivo e del quale solvoleremo la dimostrazione.

\begin{theorem}[Teorema di Lax-Richtmyer (o Teorema di equivalenza)]
    Se un metodo numerico di un problema stabile è consistente, allora è convergente:
    \[
        \textcolor{blue}{\text{consistenza} \leftrightarrow \text{convergenza}}
    \]
\end{theorem}